\documentclass[10pt]{article}
\author{Joshua N. Grant}
\title{CV for Joshua N. Grant}

\usepackage{etoolbox}
\usepackage[hidelinks]{hyperref}
\usepackage[left=1cm,top=1cm,right=1cm,bottom=1cm,nohead,nofoot]{geometry}
\usepackage{tabularx}
\usepackage{enumitem}
\usepackage{fontawesome}
\usepackage{palatino}
\usepackage{xcolor}
\usepackage[
    backend=biber,
    sorting=ydnt,
    defernumbers=true
]{biblatex}

\hypersetup{colorlinks=false}
\hypersetup{pageanchor=true}
\addbibresource{bibliography.bib}

\newcolumntype{Y}{>{\raggedright\arraybackslash}X}
\setlist[itemize]{leftmargin=1.4em,itemsep=0.1em,topsep=0.2em,parsep=0em}
\setlength{\parindent}{0pt}
\setlength{\parskip}{0.2em}

\newcommand{\sectionline}{\vspace{0.25em}\noindent\rule{\textwidth}{0.4pt}\vspace{-0.45em}}
\newcommand{\cvsection}[2]{\sectionline\section*{{\mdseries #1}\hspace{0.35em}#2}\vspace{-0.45em}}
\newcommand{\role}[4]{%
    \textbf{#1} \hfill #4\\
    \textit{#2} \hfill #3
}
\newcommand{\skillrow}[2]{\textbf{#1} & #2\\}
\newcommand{\project}[3]{%
    \textbf{#1} --- #2\\
    \textit{Technologies:} #3
}
\newcommand{\contactitem}[3]{#1\hspace{0.35em}\href{#2}{#3}}

\begin{document}
\pagestyle{empty}
\setcounter{page}{1}
\renewcommand{\thepage}{\arabic{page}}

\begin{minipage}[t]{0.43\linewidth}
    {\Huge \textbf{Joshua N. Grant}}\\[-0.1em]
    {\large Geospatial Systems Architect}\\
    {\small Geospatial Platforms \textbar{} Data Engineering \textbar{} Research Software}
\end{minipage}
\hfill
\begin{minipage}[t]{0.55\linewidth}
    \raggedleft
    {\small
    \faMapMarker\hspace{0.35em}Knoxville, TN \quad
    \faPhone\hspace{0.35em}+1 (865) 803-3495\\
    \contactitem{\faEnvelope}{mailto:jngrant@live.com}{jngrant@live.com} \quad
    \contactitem{\faGlobe}{https://sempervent.github.io}{sempervent.github.io}\\
    \contactitem{\faGithub}{https://github.com/sempervent}{github.com/sempervent} \quad
    \contactitem{\faLinkedin}{https://linkedin.com/in/joshuanagrant}{linkedin.com/in/joshuanagrant}\\
    \contactitem{\faGraduationCap}{https://scholar.google.com/citations?user=vs-HJQcAAAAJ&hl=en}{Google Scholar}
    }
\end{minipage}

\cvsection{\faUser}{Professional Summary}
Geospatial systems architect focused on turning scientific, operational, and analytic workflows into production software platforms. Designs cloud-native spatial data systems, orchestration pipelines, APIs, and decision-support applications that connect geospatial storage, processing, and delivery across backend, data, and user-facing layers. Combines geospatial engineering, data platform architecture, and research software leadership to operationalize complex data products for mission-driven environments.

\cvsection{\faCogs}{Technical Skills}
\begin{tabularx}{\linewidth}{@{}p{0.23\linewidth}Y@{}}
    \skillrow{Geospatial Systems}{PostGIS, GeoPandas, GeoParquet, GDAL, Rasterio, Pyproj, QGIS, ArcGIS, Mapbox, Leaflet, Martin, zonal statistics, raster/vector processing}
    \skillrow{Data Engineering}{Airflow, Prefect, Kafka, ETL/ELT pipelines, data modeling, workflow orchestration, Spark, Slurm, Hadoop, NiFi}
    \skillrow{Cloud \& Infrastructure}{AWS S3, Lambda, EC2, RDS, Fargate; GCP; Azure; Docker; Kubernetes; Ansible; GitHub Actions; AWS CodePipeline}
    \skillrow{Databases}{PostgreSQL, TimescaleDB, MongoDB, MySQL, Snowflake, SQLite, SpatiaLite, DuckDB, Redis, RabbitMQ}
    \skillrow{Applications \& APIs}{FastAPI, Flask, Django, REST API development, Swagger/OpenAPI, Shiny, Electron}
    \skillrow{Analytics, AI \& Visualization}{Scikit-learn, TensorFlow, Pandas, NumPy, LLM/RAG data discovery, OCR/document processing, matplotlib, Plotly, Grafana, Kibana, model deployment}
    \skillrow{Languages}{Python, SQL, R, JavaScript, Bash, PHP, \LaTeX{}}
    \skillrow{Leadership}{Technical architecture, team leadership, software quality assurance, stakeholder collaboration, Agile delivery, documentation}
\end{tabularx}

\cvsection{\faCalendar}{Professional Experience}
\role{Geospatial Systems Architect}{Oak Ridge National Laboratory}{Oak Ridge, TN}{2023--Present}
\begin{itemize}
    \item Architect geospatial data platforms that unify GeoParquet-based storage, retrieval, analysis, and downstream application delivery for mission-driven risk and decision-support workflows.
    \item Lead technical design across backend services, data models, APIs, and user-facing applications for geospatial systems deployed in both web and desktop environments.
    \item Build cloud-oriented raster and vector processing pipelines with Airflow, Prefect, Python, object storage, and geospatial libraries to support repeatable ingestion and operational analytics.
    \item Design geospatial databases and processing workflows for large raster collections, vector aggregation, zonal statistics, and clock-drift time-series analysis at operational scale.
    \item Develop high-sensitivity parcel tracking and monitoring applications using IoT streams, Kafka, TimescaleDB, Airflow, and AWS-hosted services.
    \item Improve geospatial data discovery with LLM/RAG workflows that connect users to spatial datasets, metadata, and analytic context.
    \item Modernize legacy Python systems by replacing brittle custom implementations with standard libraries, clearer interfaces, and maintainable project structure.
\end{itemize}

\role{Senior Software Engineer -- Data Engineering}{Bold Penguin}{Columbus, OH / Remote}{2022--2023}
\begin{itemize}
    \item Designed centralized development and deployment workflows using GitHub Actions, AWS services, reusable Python modules, and automated release practices, improving consistency across engineering teams.
    \item Led migration of deployments, flows, and tasks from Prefect 1 to Prefect 2, standardizing orchestration patterns and improving maintainability across teams.
    \item Built execution patterns that routed Prefect workloads to EC2 or Fargate based on runtime conditions such as file size and processing requirements.
    \item Created EC2- and Fargate-backed microservices for OCR, document conversion, data extraction, and backend processing workflows.
    \item Designed and maintained REST APIs and message queue architecture for high-load document-processing services, improving resilience and throughput in production data exchange.
    \item Automated build, test, and deployment workflows for Python and JavaScript services using GitHub Actions and AWS CodePipeline.
\end{itemize}

\role{Data Engineer}{Oak Ridge National Laboratory}{Oak Ridge, TN}{2019--2022}
\begin{itemize}
    \item Replaced a brittle cron-based web scraping process with a more reliable Airflow ETL workflow deployed in Kubernetes and cloud infrastructure.
    \item Developed full-stack Docker images and automated spatial-temporal data acquisition workflows using Python, R, Airflow, and Shiny for WSTAMP and Hadoop-backed data storage.
    \item Led development for a COVID-19 tracking project spanning R, Python, geospatial analysis, downstream model creation, and operational reporting workflows; received distinguished recognition for the project.
    \item Built containerized Airflow ETL workflows, Python microservices, and Swagger/OpenAPI REST APIs for NAERM-RTSA and related data exchange with PostgreSQL, AWS S3, AWS RDS, and MongoDB.
    \item Automated web scraping and social media ingestion for misinformation classification, statistical modeling, network analysis, Elasticsearch indexing, and Kibana visualization.
    \item Served as Software Quality Assurance board chair for the Geospatial Science and Human Security Division and contributed to ORNL Data Council discussions on data management, governance, and architecture.
\end{itemize}

\role{Data Scientist II Subcontractor}{Oak Ridge National Laboratory}{Oak Ridge, TN}{2018--2019}
\begin{itemize}
    \item Automated WSTAMP web application data ingestion using R, Python, Airflow, and Docker.
    \item Authored an R package to normalize geographic identifiers and unify data access across online APIs.
    \item Developed a SpatiaLite-backed SQLite database to simplify lookup, retrieval, and integration of geographic data.
    \item Built a Python zonal statistics application for calculating raster summaries across shapefile geometries.
\end{itemize}

\role{Owner \& Contractor-For-Hire}{Specrabella}{Knoxville, TN}{2018}
\begin{itemize}
    \item Designed and deployed interactive energy informatics applications using Leaflet, Shiny, Snowflake, and machine-learning-derived calculations.
    \item Consulted on greenfield web application architecture, database design, visualization, and deployment for energy-sector clients.
    \item Built data ingestion and visualization workflows using cURL, PHP, R, Python, and Snowflake for energy datasets.
    \item Designed PostgreSQL databases and development workflows for energy and financial data products.
\end{itemize}

\role{Bioinformatician, Project Manager, \& Technical Sales Representative}{Microbial Insights}{Knoxville, TN}{2016--2018}
\begin{itemize}
    \item Established an automated NGS analysis pipeline using R, MySQL, Python, Bash, JavaScript, and Shiny-backed internal tooling.
    \item Automated client-facing statistical workflows including linear models, ANOVA, self-organizing maps, clustering, PCA, and PCoA.
    \item Developed ETL and customer data visualization tooling using R, PHP, jQuery, and MySQL for dynamic qPCR result analysis.
    \item Built a custom R package, local Shiny applications, and an internal PostgreSQL/PHP wiki to improve reporting, project management, and customer service workflows.
\end{itemize}

\role{Graduate Research Assistant}{Dr. Neal Stewart's Plant Biotechnology Lab, University of Tennessee}{Knoxville, TN}{2014--2016}
\begin{itemize}
    \item Analyzed cell suspension experiments using linear models, ANOVA, PCA, Python, and R for DOE reporting and scientific communication.
    \item Evaluated suspension cultures through chemical and spectral processes and translated findings into research deliverables.
    \item Collaborated on a novel single-cell suspension and cryopreservation robotic system while executing monocot genetic modification workflows.
\end{itemize}

\role{Laboratory Assistant}{Dr. Neal Stewart's Plant Biotechnology Lab, University of Tennessee}{Knoxville, TN}{2012--2014}
\begin{itemize}
    \item Developed Python-based statistical methodology for screening lignin content and imaging cell characteristics both \textit{in vivo} and \textit{in vitro}.
    \item Extracted genomes from NCBI, cleaned and normalized datasets, and performed exploratory data analysis using Python.
\end{itemize}

\role{Laboratory Assistant}{Dr. Paris Lambdin's Biosystematics and Biological Control Lab, University of Tennessee}{Knoxville, TN}{2002--2012}
\begin{itemize}
    \item Designed instructional modules for undergraduate and graduate courses using HTML, CSS, JavaScript, and Flash.
    \item Performed field data collection and forest coverage analysis using ArcGIS in support of biological control and biosystematics projects.
\end{itemize}

\cvsection{\faPuzzlePiece}{Selected Projects}
\project{Where I've Been}{Travel history visualization application with spatial data storage, geospatial retrieval, and web mapping.}{Python, Flask, Docker, PostgreSQL, Leaflet, Vue, Bootstrap}

\project{Final Fantasy Football}{Predictive analytics and semantic learning application for fantasy football data exploration and decision support.}{Python, Scikit-learn, Pandas, NumPy, Flask, Docker, Airflow}

\project{Decentralized Content Reward System}{Prototype web platform exploring decentralized content incentives, workflow design, and full-stack application architecture.}{Python, Rust, Flask, Docker, PostgreSQL, React, Bootstrap}

\cvsection{\faUniversity}{Education}
\textbf{Master of Science in Plant Sciences -- Plant Molecular Genetics}\\
University of Tennessee \hfill Spring 2017\\
GPA: 3.72/4.0

\textbf{Bachelor of Science in Plant Sciences -- Biotechnology}\\
University of Tennessee \hfill Spring 2014\\
GPA: 3.74/4.0; Magna Cum Laude

\cvsection{\faBook}{Selected Publications, Technical Reports, \& Research Outputs}
\nocite{*}
\printbibliography[heading=none]

\end{document}
